\section{COMPONENTS}

\subsection{MAPBOARD}

The Mapboard shows the parts of Virginia and West Virginia in and immediately adjacent to the Shenandoah Valley in which important Civil War campaigns took place in 1862 and 1864. A hexagonal grid has been superimposed on the map to facilitate movement. Each individual hexagon is referred to as a "hex". Terrain has been somewhat simplified and, in some cases, slightly reconfigured to conform to the hexagon grid-coordinates to help pinpoint locations.

\textit{For example: Winchester is N-20, Luray is S-36, Charlottesville is BB-53, and so on.}

Partial hexes along the side of the map with 3 sides can be used for play; all other partial hexes are considered to be "off the board" and cannot be used for play.

The Mapboard also includes other features to aid in the play of SHENANDOAH. Victory Points for boths sides are tabulated on the POINT RECORD. The TURN RECORD CHART is used to keep track of the time. Two sets of Stacking Boxes are also included, one set for each player.

An arrow on the Mapboard (near the logo) shows the direction that is true North. However, for game purposes, the directions are defined differently for simplicity. North is considered to be that edge of the Mapboard that contains the grid coordinate numbers. West is the edge containing the grid coordinate numbers 1-44. South is the edge containing grid coordinate numbers 45-58. East is the remaining edge of the Mapboard.